\documentclass{jsarticle}
\usepackage[dvipdfmx, final]{graphicx}
\usepackage{amsmath}
\usepackage{amssymb}
\begin{document}
\title{仲人としてのZeta関数が
広中平祐教授とグリーシャ教授の\\
抽象理論と具象理論としての\\
ヒルベルト空間とサーストン多様体\\ 
複素多様体が統一場理論を量子力学における\\
ガウスの曲面論として\\
成し遂げる理論の背景場となるレポート}
\author{Masaaki Yamaguchi}
\date{}
\maketitle
  
  一般相対性理論における重力理論は大域的微分多様体と積分多様体についての
  単体量を求めるための空間の対数関数における不変性を記述するプランクスケール
  と異次元の宇宙におけるスケール、ウィークケールと言われているanti-D-brane
  として、この２種類の計量をガンマ関数とベータ関数としてオイラーの定数
  とそれによる連分数が微分幾何の量子化と数式の値を求めると同じという
  予知と推測値から、広中平祐定理の４重帰納法のオイラーの公式による
  多様体積分と同類値として、ペレルマン多様体がサーストン空間に成り立つ
  と同じく、広中平祐定理がヒルベルト空間で成り立つとしての、
  これら２種の多様体の場の理論が、ゼータ関数がAdS5多様体で成り立つことと、
  不変量として、ゼータ関数がこの3種の多様体の場の理論のバランスをとる
  理論として言えることである。
  　 Gravity of general relativity theory describe with Cutting of space
in being discatastrophed from Global differential and integral manifold
of scaled levetivity in plank scal and the other vector of universe
scale, these two scale inspectivity of Gamma and Beta function escourted
with these manifold experted in result of Differentail geometry of
quantum level manifold equal with Euler product and continue parameter,
moreover this product is relativity of Hironaka theorem in Four assembled
integral of Euler equation.

   \newcommand{\variint}{%
	\begin{picture}(15,30)
		\put(0,0){
			$ \iint $}
		\put(0,5){\line(15,0){15}}
	\end{picture} %
}

\newcommand{\variiint}{%
	\begin{picture}(15,15)
		\put(0,0){
			$ \iiint $}
		\put(0,5){\line(15,0){20}}
	\end{picture} %
}



 \newcommand{\alearletter}{%
	 \begin{picture}(10,20)
		 \put(0,0){
			 $ \square $}
		 \put(0,0){\line(1,1){10}}
	 \end{picture} %
	 }

 重力場理論の式は、
  Gravity equation is
  $$ \square = - {16 \pi G \over c^4}T^{\mu\nu} $$
  This equation quated with being logment of formula, and
  this formula divided with universe of number in prime zone,
  therefore, this dicided with varint equation is monotonicity of
  being composited with Weil's theorem united for Gamma function.
  この式は、対数を宇宙における数により求める素数分布論として、この
  大域的積分分断多様体がガンマ関数をヴェイル予想を根幹とする単体量
  として決まることに起因する商代数として導かれる。
  $$ \scalebox{1}[2]{\alearletter} = 
  {{{{8 \pi G} \over {c^4}}T^{\mu\nu}}/ \log x} $$
  $$ {}^{t}\scalebox{1}[2]{\variiint}\mathrm{cohom D_{\chi}[I_m]}$$
  $$ = \oint {(px^n + qx + r)^{\nabla l}} $$
  $$ {d \over d l}L(x,y) = 2 \int ||\sin 2x||^{2}d \tau $$
  $$ {d \over d \gamma}\Gamma $$
  この関数は大域的微分多様体としてのアカシックレコードの合流地点として、
  タプルスペースを形成している。
  This function esterminate with acasic record of global differential
  manifolds.
  $$ = [i \pi (\chi, x), f(x)] $$
  それにより、この多様体は基本群をアカシックレコードの相対性としての
  存在論の実存主義から統合される多様体自身としてのタプルスペースの
  池になっている。
  And this manifold from fundemental group esterminate with
  also this manifold estimate relativity of acasic record.
  $$ {d \over d \gamma}\Gamma = \int C dx_m = \int(\int {1 \over x^s}dx
  - \log x)\mathrm{dvol} $$
  また、このアカシックレコードはオイラーの定数のラムダドライバー
  にもなっている。
  More also this record tupled with lake of Euler product.
  $$ \bigoplus({i \hbar^{\nabla}})^{\oplus L} $$
  $$ x^{{1 \over 2} + iy} = e^{x \log x} $$
  それゆえに、この定数はゼータ関数と微分幾何の量子化を因数にもつ
  素因子分解の式にもなっている。
  Therefore, this product also constructed with differential geometry
  of quantum level equation and zeta function.
  $$ = C $$ 
  そして、この関数はオイラーの定数から広中平祐定理による4重帰納法の
  オイラーの公式からの多様体積分へとのサーストン空間のスペクトラム関数
  ともなっている。
  And this function of Euler product respectrum of focus with
  Heisuke Hironaka manifold in four assembled of integral Euler equation.. 

  $$ C =  b_0 + \cfrac{c_1}{b_1 +
	  \cfrac{c_2}{b_2 +
	  \cfrac{c_3}{b_3 +
	  \cfrac{c_4}{b_4 + \cdots}}}} $$
  この方程式は指数による連分数としての役割も担っている。  
  This equation demanded with continued number of step function.

  $$ {d \over d f}F(v_{ij},h) = \int {e^{f}\mathrm{dV}}
  [- \Delta v + \nabla_{i}\nabla_{j}
  v_{ij} - R_{ij}v_{ij} - v_{ij}\nabla_{i}\nabla_{j} + 2 <\nabla f,\nabla h>
  + (R + \nabla f^2)({v \over 2} - h)]$$
  $$ = {d \over d f}F = {2 {\int (R + \nabla_{i}\nabla_{j} f)^2} \over 
  { - (R + \Delta f)}}\mathrm{dm} $$
  これらの方程式は8種類の微分幾何の次元多様体として、そして、これらの
  多様体は曲平面による双対性をも生成している。 
  そして、このガウスの曲平面は、大域的微分多様体と微分幾何の量子化から
  素因数を形成してもいる。
  These equation are eight differential geometry of dimension calyement.
  And these calyement equation excluded into pair of dimension surface.
  This surface of Gauss function are global differential manifold,
  and differential geometry of quantum level.
  $$ F \ge {d \over d f}\int \int{1 \over (x \log x)^2}dx_m 
  + {d \over d f}\int \int{1 \over (y \log y)^{1 \over 2}}dy_m $$

  微分幾何の量子化はオイラーの定数とガンマ関数が指数による連分数
  としての不変性として素因数を形成していて、このガウスの曲平面
  による量子力学における重力場理論は、ダランベルシアンの切断多様体
  がこの大域的切断多様体を付加してもいる。
  Differential geometry of quantum level constructed with Euler product
  and Gamma function being discatastrophed from continued fraction style.
  Gravity equation lend with varint of monotonicity of level expresented
  from gravity of letter varient formula.
  これらの方程式は基本群と大域的微分多様体をエスコートしていもいて、
  ヴェイル予想がこのダランベルシアンの切断方程式たちから
  輸送のポートにもなっている。ベータ関数とガンマ関数がこれらのフォームラ
  の方程式を放出してもいて、結果、これらの方程式は広中平祐定理の複素
  多様体とグリーシャ教授によるペレルマン多様体からサポートされてもいる。
  この2名の教授は、一つは抽象理論をもう一方は具象理論を説明としている。
  These equation escourted into Global differential manifold and
  fundemental equation.
  Weil's theorem is imported from this equation in gravity of letters.
  Beta function and Gamma function are excluded with these formulas.
  These equation comontius from Heisuke Hironaga of complex manifold
  and Gresha professor of Perelman manifold.
  These two professos are one of abstract theorem and the other of
  visual manifold theorem.





\end{document}
